% !TeX root = ../../main.tex
\subsection{Finite-core enumeration and contributions}
\label{subsec:a:04}

We use the cores of Subsection~\ref{subsec:s03:04}: root self-loops
have been contracted, increasing the root genus. The distinguished
root is labelled $0$; the small vertices are labelled $1,\ldots,v$
temporarily for describing their edge matrix. Let $m_{ij}$ be the edge
multiplicity for $i\ne j$, and $m_{ii}$ the number of self-loops at
vertex $i$. These two-index multiplicities are distinct from the
low-degree counts $m_a$ used above. If the small genera are $G_i$,
then
\begin{equation}\label{eq:a:core-data}
\begin{gathered}
 n_i=2m_{ii}+\sum_{j\ne i}m_{ij},\qquad
 q=\sum_{i=1}^v(2G_i-2+n_i),\\
 c=\sum_{i=1}^vG_i+b_1(\Gamma_{\rm c})=\frac{q+n_0}{2}.
\end{gathered}
\end{equation}
The root has genus $g-c$ in the core. Restoring $p$ root loops gives
type $(g-c-p,n_0+2p)$ and automorphism factor $2^pp!$ in addition
to that of the core.

The enumeration is short because each small vertex has positive
complexity. We prove its completeness and the symmetry rule before
listing any contributions.

\begin{lmm}[The cores of complexity at most two]\label{lmm:a:core-enumeration}
There are respectively $1,3,13$ core types at $q=0,1,2$. For every
core their automorphism order is
\begin{equation}\label{eq:a:core-aut}
 |\operatorname{Aut}\Gamma_{\rm c}|=
 |\operatorname{Aut}_V\Gamma_{\rm c}|
       \prod_{i<j}m_{ij}!\prod_{i=0}^v2^{m_{ii}}m_{ii}!,
\end{equation}
where $\operatorname{Aut}_V$ fixes the root and consists of the
permutations of small vertices preserving genera and the multiplicity
matrix. Here $m_{00}=0$.
\end{lmm}
\begin{prf}
For fixed vertices, parallel edges may be permuted freely. A family
of $m_{ii}$ self-loops has $m_{ii}!$ permutations and two branch
exchanges per loop. A compatible vertex permutation lifts to a flag
bijection in precisely the same number of ways. This proves
\eqref{eq:a:core-aut}, including exchanges of identical small vertices.

At $q=0$ no small vertex exists. At $q=1$ the sole small vertex has
type $(0,3)$ or $(1,1)$. The equation $m_{01}+2m_{11}=n_1$,
with $m_{01}\geq1$, gives two possibilities for the former and one
for the latter.

At $q=2$ there is one complexity-two vertex or two complexity-one
vertices. A single $(0,4)$ vertex gives two possibilities and a single
$(1,2)$ gives one; type $(2,0)$ cannot connect to the root. For two
vertices put $a=m_{01}$, $b=m_{02}$, $d=m_{12}$, and $l_i=m_{ii}$.
The equations are
\[
 a+2l_1+d=n_1,\qquad b+2l_2+d=n_2.
\]
For two rational trivalent vertices, $d=0,1,2$ gives respectively
three, two, and one possibilities up to vertex exchange; $d=3$
disconnects the root. One rational trivalent and one elliptic univalent
vertex give two possibilities at $d=0$ and one at $d=1$. Two elliptic
univalent vertices have just their two separate root edges. Hence
$3+6+3+1=13$, and positivity of the small complexities excludes any
further vertex. The lists below give all the solutions.
\end{prf}

In the diagrams a solid square is the root, and a circle contains its
small vertex genus. For two small vertices the upper one is $1$ and
the lower one is $2$; for one small vertex it is $1$. All omitted
edge multiplicities are zero. A superscript $2$ on a small type denotes
two vertices of that type. The data column is
$(n_0,c;|\operatorname{Aut}\Gamma_{\rm c}|)$.

\begingroup
\newcommand{\aOne}[2]{\begin{tikzpicture}[x=1.05cm,y=.7cm,baseline=-.5ex]
 \node[rectangle,fill,minimum size=3.2mm,inner sep=0pt] (v0) at (0,0) {};
 \node[circle,draw,inner sep=1.2pt] (v1) at (1,0) {$#1$};#2
 \end{tikzpicture}}
\newcommand{\aTwo}[3]{\begin{tikzpicture}[x=1.05cm,y=.7cm,baseline=-.5ex]
 \node[rectangle,fill,minimum size=3.2mm,inner sep=0pt] (v0) at (0,0) {};
 \node[circle,draw,inner sep=1.2pt] (v1) at (1,.65) {$#1$};
 \node[circle,draw,inner sep=1.2pt] (v2) at (1,-.65) {$#2$};#3
 \end{tikzpicture}}
\begin{center}
\begin{tabular}{cclc}
 Core&Diagram&Small types and nonzero multiplicities&Data\\\hline
 $0.1$&\tikz[baseline=-.5ex]\node[rectangle,fill,minimum size=3.2mm]{};
   &none&$(0,0;1)$\\[2mm]
 $1.1$&\aOne{0}{\draw (v0)--(v1);\draw (v1) edge[loop right] (v1);}
   &$(0,3)$: $m_{01}=m_{11}=1$&$(1,1;2)$\\[2mm]
 $1.2$&\aOne{0}{\draw (v0)--(v1);\draw (v0) to[bend left=30] (v1);\draw (v0) to[bend right=30] (v1);}
   &$(0,3)$: $m_{01}=3$&$(3,2;6)$\\[2mm]
 $1.3$&\aOne{1}{\draw (v0)--(v1);}
   &$(1,1)$: $m_{01}=1$&$(1,1;1)$
\end{tabular}
\end{center}

\begin{center}
\begin{tabular}{cclc}
 Core&Diagram&Small types and nonzero multiplicities&Data\\\hline
 $2.1$&\aOne{0}{\draw (v0) to[bend left=22] (v1);\draw (v0) to[bend right=22] (v1);\draw (v1) edge[loop right] (v1);}
  &$(0,4)$: $m_{01}=2,\ m_{11}=1$&$(2,2;4)$\\[2mm]
 $2.2$&\aOne{0}{\draw (v0) to[bend left=12] (v1);\draw (v0) to[bend right=12] (v1);\draw (v0) to[bend left=40] (v1);\draw (v0) to[bend right=40] (v1);}
  &$(0,4)$: $m_{01}=4$&$(4,3;24)$\\[2mm]
 $2.3$&\aOne{1}{\draw (v0) to[bend left=22] (v1);\draw (v0) to[bend right=22] (v1);}
  &$(1,2)$: $m_{01}=2$&$(2,2;2)$\\[2mm]
 $2.4$&\aTwo{0}{0}{\draw (v0) to[bend left=18] (v1);\draw (v0) to[bend right=18] (v1);\draw (v1)--(v2);\draw (v2) edge[loop right] (v2);}
  &$(0,3)^2$: $m_{01}=2,\ m_{12}=m_{22}=1$&$(2,2;4)$\\[2mm]
 $2.5$&\aTwo{0}{0}{\draw (v0)--(v1);\draw (v0)--(v2);\draw (v1) to[bend left=20] (v2);\draw (v1) to[bend right=20] (v2);}
  &$(0,3)^2$: $m_{01}=m_{02}=1,\ m_{12}=2$&$(2,2;4)$\\[2mm]
 $2.6$&\aTwo{0}{0}{\draw (v0)--(v1);\draw (v0)--(v2);\draw (v1) edge[loop right] (v1);\draw (v2) edge[loop right] (v2);}
  &$(0,3)^2$: $m_{01}=m_{02}=m_{11}=m_{22}=1$&$(2,2;8)$
\end{tabular}
\end{center}

\begin{center}
\begin{tabular}{cclc}
 Core&Diagram&Small types and nonzero multiplicities&Data\\\hline
 $2.7$&\aTwo{0}{0}{\draw (v0)--(v1);\draw (v0) to[bend left=28] (v1);\draw (v0) to[bend right=28] (v1);\draw (v0)--(v2);\draw (v2) edge[loop right] (v2);}
  &$(0,3)^2$: $m_{01}=3,\ m_{02}=m_{22}=1$&$(4,3;12)$\\[2mm]
 $2.8$&\aTwo{0}{0}{\draw (v0) to[bend left=18] (v1);\draw (v0) to[bend right=18] (v1);\draw (v0) to[bend left=18] (v2);\draw (v0) to[bend right=18] (v2);\draw (v1)--(v2);}
  &$(0,3)^2$: $m_{01}=m_{02}=2,\ m_{12}=1$&$(4,3;8)$\\[2mm]
 $2.9$&\aTwo{0}{0}{\draw (v0)--(v1);\draw (v0) to[bend left=28] (v1);\draw (v0) to[bend right=28] (v1);\draw (v0)--(v2);\draw (v0) to[bend left=28] (v2);\draw (v0) to[bend right=28] (v2);}
  &$(0,3)^2$: $m_{01}=m_{02}=3$&$(6,4;72)$\\[2mm]
 $2.10$&\aTwo{0}{1}{\draw (v0)--(v1);\draw (v0)--(v2);\draw (v1) edge[loop right] (v1);}
  &$(0,3),(1,1)$: $m_{01}=m_{02}=m_{11}=1$&$(2,2;2)$\\[2mm]
 $2.11$&\aTwo{0}{1}{\draw (v0) to[bend left=18] (v1);\draw (v0) to[bend right=18] (v1);\draw (v1)--(v2);}
  &$(0,3),(1,1)$: $m_{01}=2,\ m_{12}=1$&$(2,2;2)$\\[2mm]
 $2.12$&\aTwo{0}{1}{\draw (v0)--(v1);\draw (v0) to[bend left=28] (v1);\draw (v0) to[bend right=28] (v1);\draw (v0)--(v2);}
  &$(0,3),(1,1)$: $m_{01}=3,\ m_{02}=1$&$(4,3;6)$\\[2mm]
 $2.13$&\aTwo{1}{1}{\draw (v0)--(v1);\draw (v0)--(v2);}
  &$(1,1)^2$: $m_{01}=m_{02}=1$&$(2,2;2)$
\end{tabular}
\end{center}
\endgroup

For example the two rational vertices of core $2.9$ can be exchanged,
in addition to permuting both triples of parallel edges: its order is
$2(3!)^2=72$. Core $2.4$ has no vertex exchange and has order
$2!\cdot2=4$. Both elliptic vertices of core $2.13$ can be exchanged,
giving its order two. These distinctions matter even though no
markings occur in these graphs.

We next give the edge integration rule and its complete evaluation.
For a complex upper endpoint $z$, integration is along the straight
segment from zero to $z$. Put, for $a=0,1,2$,
\begin{equation}\label{eq:a:core-edge-moments}
 \mathcal J_a(z)=\int_0^z t^a f(t)\,\dd t,\qquad
 E_f(z)=\frac{\pi^2\mathcal J_0(z)}{12}
                         -\frac{\mathcal J_1(z)}{24},\qquad
 E_0(z)=\frac{\pi^2z}{12}-\frac{z^2}{48}.
\end{equation}
These are entire functions. The two $E$ functions integrate the
elliptic univalent vertex on a root edge, with its factor $f$, or on
an edge between small vertices, respectively. They and $\mathcal J_a$
are used only to state the following core integrals.

Define $\mathcal I_{\Gamma_{\rm c}}(z)$ by attaching $f(t_e)$ to
every root flag, the exact polynomial~\eqref{eq:a:small-vertices}
to every small vertex, integrating $t_e$ over $[0,z]$ for every core
edge, and dividing by $|\operatorname{Aut}\Gamma_{\rm c}|$. A small
self-loop supplies the same $t_e$ twice to its vertex polynomial.
There is no additional vertex-label factorial.

\begin{lmm}[Leading core integration]\label{lmm:a:core-integration}
For a fixed core of complexity $q$, the leading contribution to
$Z_g(u)/V_{g,0}$ is
\begin{equation}\label{eq:a:leading-core}
 \frac{e^{S(u)}}{(8\pi^2g)^q}
           \mathcal I_{\Gamma_{\rm c}}(\alpha u).
\end{equation}
For the absolute normalization multiply by $1/\sqrt\pi$. The error
in either normalization is $O(g^{-q-1})$, locally uniformly in complex $u$.
\end{lmm}
\begin{prf}
The exact core volume path~\eqref{eq:s03:core-volume-path} has
leading scale $(8\pi^2g)^{-q}$, independently of the restored loop
count $p$. The leading root profile is a product of $f$ over its
flags. Integrating the $p$ restored root loops, and using their factor
$2^pp!$, gives
\[
 \frac1{p!}\left(\frac12\int_0^{\alpha u}f(t)^2\dd t\right)^p
 =\frac{S(u)^p}{p!}.
\]
The remaining integrals are precisely $\mathcal I_{\Gamma_{\rm c}}$.
The summable core remainder from Section~\ref{sec:s03} proves that
the sum over all legal $p$ has the asserted error, and permits extending
it to the entire exponential series. Finally use
$V_{g,0}/F_g=\pi^{-1/2}(1+O(g^{-1}))$.
\end{prf}

The endpoint evaluations needed in all rows are especially simple:
\begin{equation}\label{eq:a:endpoint-moments}
\begin{gathered}
 z=\alpha=2\pi^2,\qquad \mathcal J_0=8,\qquad
 \mathcal J_1=8\pi^2-32,\\
 \mathcal J_2=16\pi^4-192\pi^2+768,\qquad
 E_f=(\pi^2+4)/3,\qquad E_0=\pi^4/12.
\end{gathered}
\end{equation}
Indeed $t=2y^2$ gives
$\mathcal J_a(\alpha)=2^{a+2}\int_0^\pi y^{2a}\sin y\dd y$.
The latter integral is $2$ at $a=0$, and successive integration by
parts subtracts $(2a)(2a-1)$ times the preceding integral from
$\pi^{2a}$. This proves all the displayed values.

In the tables below $\mathcal J_a,E_f,E_0$ are evaluated at the
variable $z$ in the integral column and at $z=\alpha$ in the last
column. The last column is always
$\mathcal I_{\Gamma_{\rm c}}(\alpha)/(8\pi^2)^q$. Thus it is the
coefficient of $e^{S(1)}g^{-q}$ in the relative normalization, or of
$e^{S(1)}g^{-q}/\sqrt\pi$ in the absolute normalization.

\begin{center}
\begin{tabular}{ccl}
 Core&$\mathcal I_{\Gamma_{\rm c}}(z)$&Normalized endpoint\\[1mm]\hline
 \noalign{\smallskip}
 $0.1$&$1$&$1$\\[1mm]
 $1.1$&$z\mathcal J_0/2$&$1$\\[1mm]
 $1.2$&$\mathcal J_0^3/6$&$32/(3\pi^2)$\\[1mm]
 $1.3$&$E_f$&$1/24+1/(6\pi^2)$
\end{tabular}
\end{center}

The three rows with a single complexity-two vertex are
\begin{center}
\begin{tabular}{ccl}
 Core&$\mathcal I_{\Gamma_{\rm c}}(z)$&Normalized endpoint\\[1mm]\hline
 \noalign{\smallskip}
 $2.1$&$\bigl[(2\pi^2z-z^2)\mathcal J_0^2
                              -2z\mathcal J_0\mathcal J_1\bigr]/4$
 & $-1+4/\pi^2$\\[2mm]
 $2.2$&$\bigl[2\pi^2\mathcal J_0^4
                         -4\mathcal J_0^3\mathcal J_1\bigr]/24$
 & $-16/(3\pi^2)+128/(3\pi^4)$\\[2mm]
 $2.3$&$\begin{gathered}
 \pi^4\mathcal J_0^2/8-\pi^2\mathcal J_0\mathcal J_1/6\\
              +(\mathcal J_0\mathcal J_2+\mathcal J_1^2)/48
 \end{gathered}$
 & $1/48+7/(3\pi^4)$
\end{tabular}
\end{center}

The six rows with two rational trivalent vertices are
\begin{center}
\begin{tabular}{ccl}
 Core&$\mathcal I_{\Gamma_{\rm c}}(z)$&Normalized endpoint\\[1mm]\hline
 \noalign{\smallskip}
 $2.4$&$z^2\mathcal J_0^2/4$&$1$\\[1mm]
 $2.5$&$z^2\mathcal J_0^2/4$&$1$\\[1mm]
 $2.6$&$z^2\mathcal J_0^2/8$&$1/2$\\[1mm]
 $2.7$&$z\mathcal J_0^4/12$&$32/(3\pi^2)$\\[1mm]
 $2.8$&$z\mathcal J_0^4/8$&$16/\pi^2$\\[1mm]
 $2.9$&$\mathcal J_0^6/72$&$512/(9\pi^4)$
\end{tabular}
\end{center}

Finally the three mixed rows and the two-elliptic row are
\begin{center}
\begin{tabular}{ccl}
 Core&$\mathcal I_{\Gamma_{\rm c}}(z)$&Normalized endpoint\\[1mm]\hline
 \noalign{\smallskip}
 $2.10$&$z\mathcal J_0E_f/2$&$1/24+1/(6\pi^2)$\\[1mm]
 $2.11$&$\mathcal J_0^2E_0/2$&$1/24$\\[1mm]
 $2.12$&$\mathcal J_0^3E_f/6$&$4/(9\pi^2)+16/(9\pi^4)$\\[1mm]
 $2.13$&$E_f^2/2$&$1/1152+1/(144\pi^2)+1/(72\pi^4)$
\end{tabular}
\end{center}

Every integral column follows by multiplying the small vertex
polynomials. For instance, in core $2.1$ the small loop has parameter
$t$ twice, so integrating the polynomial
$2\pi^2-t_1-t_2-2t$ over this loop gives
$(2\pi^2z-z^2)-z(t_1+t_2)$. The two root edges then give the first
integral in that table, with automorphism divisor four. This explains
in particular its negative contribution at $z=\alpha$.

\begin{prp}[Finite contribution summary]\label{prp:a:finite-summary}
The complexity-two leading group is
\begin{equation}\label{eq:a:q-two}
\begin{aligned}
 \left.\frac{Z_g(1)}{F_g}\right|_{q=2}
 &=\frac{e^{S(1)}}{\sqrt\pi g^2}
 \left(\frac{1849}{1152}+\frac{3737}{144\pi^2}
                             +\frac{7465}{72\pi^4}\right)\\
 &\quad+O(g^{-3}).
\end{aligned}
\end{equation}
Together with~\eqref{eq:a:q-zero} and~\eqref{eq:a:q-one}, the finite
groups have total coefficients $9/8$ at $g^{-1}$ and
$257/128+1/(8\pi^2)$ at $g^{-2}$, after extracting
$e^{S(1)}/\sqrt\pi$.
\end{prp}
\begin{prf}
The complete table, whose completeness was proved above, gives
\begin{align*}
 \sum_{\Gamma_{\rm c}:\,q=1}\mathcal I_{\Gamma_{\rm c}}(\alpha)
                 &=\frac{25\pi^2+260}{3},\\
 \sum_{\Gamma_{\rm c}:\,q=2}\mathcal I_{\Gamma_{\rm c}}(\alpha)
                 &=\frac{1849\pi^4}{18}
                         +\frac{14948\pi^2}{9}+\frac{59720}{9}.
\end{align*}
Here each sum runs over core types of the indicated complexity, once
each. Divide the second by $(8\pi^2)^2$ and apply
Lemma~\ref{lmm:a:core-integration}. The first agrees independently
with the leading column of~\eqref{eq:a:q-one-table}.
For the claimed totals, add the proved $q=0,1,2$ coefficients. Their
$\pi^{-4}$ terms cancel at second order, and their $\pi^{-2}$ terms
leave $1/8$. The global tail for $q\geq3$ and graphs without a large
vertex is supplied by~\eqref{eq:s03:graph-tail}; its application to
the Hodge normalization is made in Section~4.
\end{prf}

\begin{rmk}[Reproducibility]
The accompanying calculations use rational arithmetic in $\pi^2$
to check the profiles, moments and ratio rules. A separate computation,
using only the Python standard library, solves
the valence equations, counts literal flag automorphisms, derives the
four low vertex polynomials from string and dilaton, and integrates
each row. It does not take the counts $1,3,13$ or any contribution
coefficient as an enumeration input. The installation and commands
are documented in the calculations directory; running them writes
only to the ignored build directory or to the console.
\end{rmk}
